\documentclass[10pt,a4paper]{article}

% ---------- packages ----------
\usepackage[margin=2.5cm]{geometry}
\usepackage{listings}
\usepackage{xcolor}
\usepackage{booktabs}
\usepackage{hyperref}
\usepackage{parskip}
\usepackage{fancyhdr}
\usepackage{microtype}
\usepackage{amssymb}

% ---------- page style ----------
\pagestyle{fancy}
\fancyhf{}
\rhead{CSE341 — D3}
\lhead{WorkflowScript}
\rfoot{\thepage}

% ---------- colours ----------
\definecolor{kw}{RGB}{0,0,180}
\definecolor{cm}{RGB}{80,130,80}
\definecolor{str}{RGB}{163,21,21}
\definecolor{bg}{RGB}{248,248,248}
\definecolor{err}{RGB}{180,0,0}
\definecolor{out}{RGB}{30,100,30}
\definecolor{frame}{RGB}{200,200,200}

% ---------- WorkflowScript listing style ----------
\lstdefinelanguage{WorkflowScript}{
  morekeywords={task,pipeline,stage,on_failure,fun,var,run,return,
                if,else,while,int,float,bool,void,
                timeout,retries,parallel,true,false},
  sensitive=true,
  morecomment=[l]{//},
  literate={->}{{$\to$}}1
           {<=}{{$\leq$}}1
           {>=}{{$\geq$}}1
           {!=}{{$\neq$}}1
           {==}{{$=$}}1
}

\lstset{
  language=WorkflowScript,
  basicstyle=\ttfamily\small,
  keywordstyle=\color{kw}\bfseries,
  commentstyle=\color{cm}\itshape,
  backgroundcolor=\color{bg},
  frame=single,
  rulecolor=\color{frame},
  framesep=4pt,
  xleftmargin=12pt,
  xrightmargin=4pt,
  numbers=left,
  numberstyle=\tiny\color{gray},
  numbersep=8pt,
  breaklines=true,
  showstringspaces=false,
  tabsize=4,
  captionpos=b,
}

\lstdefinestyle{error}{
  language={},
  basicstyle=\ttfamily\small\color{err},
  backgroundcolor=\color{bg},
  frame=single,
  rulecolor=\color{frame},
  framesep=4pt,
  xleftmargin=12pt,
  numbers=none,
  breaklines=true,
}

\lstdefinestyle{output}{
  language={},
  basicstyle=\ttfamily\small\color{out},
  backgroundcolor=\color{bg},
  frame=single,
  rulecolor=\color{frame},
  framesep=4pt,
  xleftmargin=12pt,
  numbers=none,
  breaklines=true,
}

% =========================================================
\begin{document}
% =========================================================

\thispagestyle{fancy}

\bigskip
\noindent
All programs were tested with the \texttt{wsc} CLI\@.
Valid programs are executed with \texttt{wsc <file>}; interpreter output is written to
stdout.  Programs with type errors are rejected by the type checker before any execution
begins.  Malformed programs are rejected at parse time with a diagnostic message on stderr.

% =========================================================
\section{Valid Programs --- Execution with Outputs}
\label{sec:valid}
% =========================================================

The table below summarises which language features each example exercises.

\begin{center}
\small
\begin{tabular}{lccc}
\toprule
Feature & Ex.~1 & Ex.~2 & Ex.~3 \\
\midrule
\texttt{task} declaration (structured type)  & \checkmark & \checkmark & \checkmark \\
Function definition (\texttt{fun})           & \checkmark & \checkmark & \checkmark \\
Function call                                & \checkmark & \checkmark & \checkmark \\
Task field access (\texttt{t.field})         & \checkmark & \checkmark & \checkmark \\
\texttt{if} statement                        & \checkmark & \checkmark & \checkmark \\
\texttt{else} branch                         & \checkmark &            &            \\
\texttt{while} loop                          &            &            & \checkmark \\
\texttt{pipeline} + \texttt{stage}           &            &            & \checkmark \\
\texttt{on\_failure}                         &            &            & \checkmark \\
\texttt{run} statement                       & \checkmark & \checkmark & \checkmark \\
\bottomrule
\end{tabular}
\end{center}

% ---------------------------------------------------------
\subsection{Program 1 --- Conditional Task Dispatch (\texttt{01\_basic.ws})}
% ---------------------------------------------------------

\begin{lstlisting}[caption={01\_basic.ws --- conditional dispatch based on task cost}]
// Conditional dispatch: choose the cheaper task based on computed cost
task fast {
    timeout:  20
    retries:  0
    parallel: true
}

task thorough {
    timeout:  90
    retries:  3
    parallel: false
}

fun total_cost(t: task) -> int {
    return t.timeout * (t.retries + 1)
}

var fast_cost: int     = total_cost(fast)
var thorough_cost: int = total_cost(thorough)

if fast_cost < thorough_cost {
    run fast
} else {
    run thorough
}
\end{lstlisting}

\noindent\textbf{Actual output:}
\begin{lstlisting}[style=output]
[run] fast
\end{lstlisting}

\noindent\textbf{Discussion.}
This program demonstrates user-defined functions that accept a \texttt{task} argument and
access its fields (\texttt{t.timeout}, \texttt{t.retries}) to compute a derived integer
value.
The two tasks, \texttt{fast} and \texttt{thorough}, are passed to \texttt{total\_cost},
which returns $20 \times 1 = 20$ and $90 \times 4 = 360$ respectively.
Because \texttt{fast\_cost} $<$ \texttt{thorough\_cost} is true, the \texttt{if} branch
executes and \texttt{run fast} is triggered, printing \texttt{[run] fast}.
The program exercises WorkflowScript's structured type (\texttt{task}), field access via
dot notation, function calls with task parameters, arithmetic expressions, and the
\texttt{if-else} control structure in a single coherent workflow scenario.

% ---------------------------------------------------------
\subsection{Program 2 --- Function Definitions and Task Field Access (\texttt{02\_functions.ws})}
% ---------------------------------------------------------

\begin{lstlisting}[caption={02\_functions.ws --- function definitions and task field access}]
// User-defined functions and task field access
task deploy { timeout: 120 retries: 1 parallel: false }

fun max_time(t: int, r: int) -> int {
    return t * (r + 1)
}

fun is_long(t: task) -> bool {
    return t.timeout > 60
}

var m: int    = max_time(120, 1)
var slow: bool = is_long(deploy)

if slow {
    run deploy
}
\end{lstlisting}

\noindent\textbf{Actual output:}
\begin{lstlisting}[style=output]
[run] deploy
\end{lstlisting}

\noindent\textbf{Discussion.}
This program showcases two function definitions with different return types
(\texttt{int} and \texttt{bool}) to demonstrate that WorkflowScript's type system
correctly tracks return types across function calls.
\texttt{max\_time} receives two \texttt{int} arguments and returns their product,
computing the maximum possible wall-clock duration for a task with retries.
\texttt{is\_long} accepts a \texttt{task} argument and returns a boolean by comparing
\texttt{t.timeout > 60}; since \texttt{deploy.timeout} is 120, it returns \texttt{true}.
The \texttt{if} body therefore executes and \texttt{[run] deploy} is printed, illustrating
how functions that access structured type fields and produce boolean results can drive
control flow.

% ---------------------------------------------------------
\subsection{Program 3 --- Pipeline Execution with While Loop (\texttt{03\_pipeline.ws})}
% ---------------------------------------------------------

\begin{lstlisting}[caption={03\_pipeline.ws --- pipeline with on\_failure, while loop, and conditional run}]
// Full pipeline with on_failure, functions, while loop
task build  { timeout: 30  retries: 0 parallel: false }
task test   { timeout: 60  retries: 2 parallel: true  }
task deploy { timeout: 120 retries: 1 parallel: false }

pipeline release {
    stage build
    stage test
    stage deploy
    on_failure { run build }
}

fun should_run(limit: int) -> bool {
    return limit > 0
}

var attempt: int = 0
var limit: int   = 3

while attempt < limit {
    attempt = attempt + 1
}

if should_run(limit) {
    run release
}
\end{lstlisting}

\noindent\textbf{Actual output:}
\begin{lstlisting}[style=output]
[run] build
[run] test
[run] deploy
\end{lstlisting}

\noindent\textbf{Discussion.}
This program is the most comprehensive of the three: it declares three tasks, a named
pipeline with three ordered stages and an \texttt{on\_failure} handler, a \texttt{while}
loop that counts from 0 to 3, and a guard function that decides whether the pipeline
should run.
The \texttt{while} loop exits when \texttt{attempt == limit == 3}; then
\texttt{should\_run(3)} returns \texttt{true}, so \texttt{run release} is executed.
Running a pipeline sequentially prints each stage name in declaration order ---
\texttt{[run] build}, \texttt{[run] test}, \texttt{[run] deploy} --- which matches
the \textsc{[Run-Pipeline]} semantic rule in the design specification.
The \texttt{on\_failure} block and \texttt{parallel: true} field on \texttt{test} are
parsed and type-checked but are not triggered during this simulated execution,
demonstrating their presence as first-class language constructs even without a real
process scheduler.

% =========================================================
\section{Type Checker --- Type Error Caught Before Execution}
\label{sec:typeerror}
% =========================================================

% ---------------------------------------------------------
\subsection{Type Error --- Bool Field Used in Arithmetic (\texttt{06\_type\_error.ws})}
% ---------------------------------------------------------

\begin{lstlisting}[caption={06\_type\_error.ws --- bool task field in arithmetic expression}]
// Type error: mixing bool task field in arithmetic expression.
// deploy.parallel is of type bool; the + operator requires numeric operands.
// The type checker catches this before any execution begins.
task deploy { timeout: 120  retries: 1  parallel: false }

var effective_cost: int = deploy.timeout + deploy.parallel
\end{lstlisting}

\noindent\textbf{Type error message:}
\begin{lstlisting}[style=error]
type error: line 6: + requires numeric operands
\end{lstlisting}

\noindent\textbf{Discussion.}
This program contains a common domain mistake: a programmer attempts to include the
\texttt{parallel} flag in a cost formula alongside \texttt{timeout}.
In WorkflowScript, \texttt{deploy.parallel} has type \texttt{bool}, while
\texttt{deploy.timeout} has type \texttt{int}; the \texttt{+} operator requires both
operands to be numeric (\texttt{int} or \texttt{float}), so applying it to a
\texttt{bool} operand is forbidden by the type rules in \S4.5.2 of the design
specification.
The type checker detects this before any execution begins and halts with a precise error
on line~6.
No output is produced because interpretation never starts --- this demonstrates
WorkflowScript's \emph{fail-fast} design principle: type safety is enforced at compile
time to prevent silent failures during long-running pipeline jobs.

% =========================================================
\section{Malformed Programs --- Parser Error Messages}
\label{sec:invalid}
% =========================================================

Each subsection gives the source file, the exact error string printed to \texttt{stderr}
by \texttt{wsc}, and a brief explanation of which grammar rule the input violates.

% ---------------------------------------------------------
\subsection{Error 1 --- Missing Closing Brace (\texttt{01\_missing\_brace.ws})}
% ---------------------------------------------------------

\begin{lstlisting}[caption={01\_missing\_brace.ws}]
task build { timeout: 30 retries: 0 parallel: false
\end{lstlisting}

\begin{lstlisting}[style=error]
line 2: expected RBRACE, got ""
\end{lstlisting}

\noindent\textbf{Explanation.}
The grammar rule for \texttt{<task\_decl>} requires a closing \texttt{\}} after the
three field declarations.  The source ends immediately after the \texttt{false} token,
so the parser reaches end-of-file while still expecting \texttt{RBRACE}.
The ``got \texttt{""}'' indicates the EOF token, which has an empty lexeme.

% ---------------------------------------------------------
\subsection{Error 2 --- Wrong Task Field Order (\texttt{02\_bad\_field\_order.ws})}
% ---------------------------------------------------------

\begin{lstlisting}[caption={02\_bad\_field\_order.ws}]
task build { retries: 0 timeout: 30 parallel: false }
\end{lstlisting}

\begin{lstlisting}[style=error]
line 1: expected KW_TIMEOUT, got "retries"
\end{lstlisting}

\noindent\textbf{Explanation.}
The grammar mandates a fixed field order: \texttt{timeout}, then \texttt{retries},
then \texttt{parallel}.  Swapping \texttt{retries} and \texttt{timeout} violates this
ordering; the parser immediately reports that it expected the keyword \texttt{timeout}
but found \texttt{retries} instead.

% ---------------------------------------------------------
\subsection{Error 3 --- Empty Pipeline (\texttt{03\_empty\_pipeline.ws})}
% ---------------------------------------------------------

\begin{lstlisting}[caption={03\_empty\_pipeline.ws}]
pipeline ci { }
\end{lstlisting}

\begin{lstlisting}[style=error]
line 1: pipeline must have at least one stage
\end{lstlisting}

\noindent\textbf{Explanation.}
The grammar rule \texttt{<pipeline\_body>~::=~<stage\_stmt>~\{~<stage\_stmt>~\}~[~<on\_failure>~]}
requires at least one \texttt{stage} statement.  An empty body violates this constraint.
The parser checks \texttt{len(stages) == 0} after the stage-collection loop and emits a
semantic parse error before the closing \texttt{\}} is consumed.

% ---------------------------------------------------------
\subsection{Error 4 --- Missing Arrow in Function Signature (\texttt{04\_missing\_arrow.ws})}
% ---------------------------------------------------------

\begin{lstlisting}[caption={04\_missing\_arrow.ws}]
fun double(x: int) int { return x }
\end{lstlisting}

\begin{lstlisting}[style=error]
line 1: expected ARROW, got "int"
\end{lstlisting}

\noindent\textbf{Explanation.}
The grammar requires the \texttt{->} separator between the parameter list and the
return type: \texttt{<fun\_decl>~::=~"fun"~IDENT~"("~\ldots~")"~"->"~<type>~<block>}.
Omitting \texttt{->} causes the parser to see the return-type keyword \texttt{int}
where it expects the \texttt{ARROW} token.

% ---------------------------------------------------------
\subsection{Error 5 --- Integer Literal for Boolean Field (\texttt{05\_bad\_parallel.ws})}
% ---------------------------------------------------------

\begin{lstlisting}[caption={05\_bad\_parallel.ws}]
task build { timeout: 30 retries: 0 parallel: 1 }
\end{lstlisting}

\begin{lstlisting}[style=error]
line 1: expected BOOL_LIT, got "1"
\end{lstlisting}

\noindent\textbf{Explanation.}
The grammar specifies \texttt{"parallel"~":"~BOOL\_LIT}; only the tokens \texttt{true}
and \texttt{false} satisfy \texttt{BOOL\_LIT}.  The integer literal \texttt{1} is
lexed as \texttt{INT\_LIT}, which does not match the expected token type.
This demonstrates that WorkflowScript does not perform implicit numeric-to-boolean
coercion, even at parse time.

% =========================================================
\end{document}
% =========================================================
